\documentclass[10.5pt,a4paper]{article}
\usepackage[utf8]{inputenc}
\usepackage[T1]{fontenc}
\usepackage{lmodern}
\usepackage[margin=2.1cm]{geometry}
\usepackage{amsmath,amssymb}
\usepackage{booktabs,longtable,array,tabularx}
\usepackage{xcolor}
\usepackage[hidelinks]{hyperref}
\usepackage{fancyhdr}

\definecolor{accent}{HTML}{1F4E79}
\definecolor{soft}{HTML}{EAF1F8}
\definecolor{warn}{HTML}{FFF4E0}
\definecolor{warnline}{HTML}{C77700}

\makeatletter
\renewcommand\section{\@startsection{section}{1}{0pt}{1.5em plus .3em minus .2em}{0.5em}{\Large\bfseries\color{accent}}}
\renewcommand\subsection{\@startsection{subsection}{2}{0pt}{1.2em plus .2em minus .1em}{0.4em}{\large\bfseries\color{accent}}}
\makeatother
\setlength{\parindent}{0pt}
\setlength{\parskip}{0.45em}
\renewcommand{\arraystretch}{1.25}
\let\olditemize\itemize\renewcommand\itemize{\olditemize\setlength\itemsep{0.15em}\setlength\topsep{0.3em}}
\let\oldenumerate\enumerate\renewcommand\enumerate{\oldenumerate\setlength\itemsep{0.15em}\setlength\topsep{0.3em}}

\pagestyle{fancy}
\fancyhf{}
\fancyhead[L]{\small Paper manual: Sanz et al., PRL 99, 055501 (2007)}
\fancyhead[R]{\small page \thepage}
\renewcommand{\headrulewidth}{0.3pt}

\newenvironment{keybox}{\par\medskip\noindent\begin{lrbox}{\mybox}\begin{minipage}{\dimexpr\linewidth-2\fboxsep-2\fboxrule\relax}}{\end{minipage}\end{lrbox}\noindent\fcolorbox{accent}{soft}{\usebox{\mybox}}\par\medskip}
\newenvironment{warnbox}{\par\medskip\noindent\begin{lrbox}{\mybox}\begin{minipage}{\dimexpr\linewidth-2\fboxsep-2\fboxrule\relax}}{\end{minipage}\end{lrbox}\noindent\fcolorbox{warnline}{warn}{\usebox{\mybox}}\par\medskip}
\newsavebox{\mybox}

\newcommand{\Ts}{T^{*}}
\newcommand{\ps}{p^{*}}
\newcommand{\kT}{k_{\mathrm B}T}
\newcommand{\yes}{\textcolor{green!45!black}{\textbf{yes}}}
\newcommand{\no}{\textcolor{red!70!black}{\textbf{no}}}
\newcommand{\done}{\textcolor{green!45!black}{\textbf{done}}}
\newcommand{\notdone}{\textcolor{red!70!black}{\textbf{not done}}}

\begin{document}

\begin{center}
{\LARGE\bfseries\color{accent} Paper manual\\[2pt]}
{\large Evidence for Out-of-Equilibrium Crystal Nucleation in Suspensions of Oppositely Charged Colloids}\\[4pt]
E.\ Sanz, C.\ Valeriani, D.\ Frenkel, M.\ Dijkstra, \emph{Phys.\ Rev.\ Lett.} \textbf{99}, 055501 (2007)\\[4pt]
\small What they did, in order, how each number was measured, and how to redo the part you reproduced.
\end{center}

\begin{keybox}
\textbf{How to use this document.} Read sections 1--4 once to get the story (about 15 minutes). Use sections 5--8 as a reference when you redo or explain something.
Every number was checked against the PDF pages. Where I read a value off a figure it is marked \emph{approx.}\ Where the paper is silent, the table in section~2 says so.
\end{keybox}

%--------------------------------------------------------------------
\section{The objective}

A liquid freezes through a rare event: a small crystal nucleus must form and grow beyond a \textbf{critical size}. When more than one crystal is possible, \emph{which one forms?}

The textbook answer is the \textbf{Stranski--Totomanow (ST) conjecture}: the crystal with the \textbf{lowest free-energy barrier} nucleates first.
In classical nucleation theory (CNT) the rate is
\begin{equation}
R=\Gamma\,e^{-\Delta G_{\mathrm{crit}}/\kT},
\end{equation}
where $\Gamma$ is a kinetic prefactor and $\Delta G_{\mathrm{crit}}$ the barrier. ST rests on two hidden assumptions: (i) sub-critical clusters are in \emph{quasi-equilibrium} (they keep rearranging into their lowest-free-energy structure), and (ii) $\Gamma$ is similar for the competing crystals, so only the barrier matters.

The paper tests ST by measuring two things \textbf{independently} and checking whether they agree.

\begin{center}
\begin{tabularx}{\textwidth}{@{}>{\bfseries}l X X@{}}
\toprule
Quantity & Question it answers & Method\\
\midrule
Nucleation \textbf{rate} & Which crystal \emph{actually} nucleates, and how fast? & Forward-flux sampling (FFS)\\
Free-energy \textbf{barrier} & Which crystal is \emph{thermodynamically cheapest} to nucleate? & Umbrella sampling\\
\bottomrule
\end{tabularx}
\end{center}

\begin{keybox}
\textbf{What would disprove ST:} the crystal that nucleates has the \emph{higher} barrier, and the nucleating structure \emph{changes} when you speed up the equilibration of the cluster with swap moves. The second observation shows that the clusters were not in equilibrium.
\end{keybox}

%--------------------------------------------------------------------
\section{The system and every input}

\textbf{Model.} A 1:1 mixture of hard spheres of diameter $\sigma$, half positive and half negative, interacting through a screened-Coulomb (Yukawa) pair potential (paper Eq.~2):
\begin{equation}
\frac{u(r)}{\kT}=
\begin{cases}
\infty & r<\sigma\\[2pt]
\pm\dfrac{Z^{2}}{(1+\kappa\sigma/2)^{2}}\,\dfrac{\lambda_B}{\sigma}\,\dfrac{e^{-\kappa(r-\sigma)}}{r/\sigma} & r_c>r\ge\sigma\\[8pt]
0 & r\ge r_c
\end{cases}
\end{equation}
The sign is $+$ for equal charges and $-$ for opposite charges. Here $\lambda_B=e^{2}/\epsilon_s\kT$ is the Bjerrum length and $\kappa=\sqrt{8\pi\lambda_B\rho_{\mathrm{salt}}}$ the inverse Debye length.

\textbf{Reduced units.}
\begin{equation}
\Ts=\frac{(1+\kappa\sigma/2)^{2}\,\sigma}{Z^{2}\lambda_B}\;=\;\frac{\kT}{\text{contact energy}},\qquad
\ps=\frac{p\,\sigma^{3}}{\kT},\qquad
\eta=\frac{\pi\sigma^{3}N}{6V}.
\end{equation}
In your code (\texttt{Source/ener.f}) the pair energy is $u=\pm e^{-\kappa(r-\sigma)}/(r/\sigma)$ in units of the contact energy, which is the same potential.

\subsection*{Input table}

\small
\begin{longtable}{@{}>{\raggedright\arraybackslash}p{4.1cm} >{\raggedright\arraybackslash}p{4.6cm} c >{\raggedright\arraybackslash}p{4.5cm}@{}}
\toprule
\textbf{Input} & \textbf{Paper value} & \textbf{Stated?} & \textbf{Your code / runs}\\
\midrule
\endhead
Composition & 1:1 binary & \yes & 500 + / 500 $-$\\
$\kappa\sigma$ & 6 & \yes & 6\\
Cutoff $r_c$ & $3.5\sigma$ & \yes & 3.5\\
$\Ts$ & 1 (disordered-fcc/CsCl/fluid triple point at 1.07) & \yes & 1\\
$\ps$ for FFS and umbrella sampling & \textbf{15} ($\ps/\ps_{\rm coex}\approx1.7$, $\ps_{\rm coex}=8.8\pm0.1$) & \yes & not used\\
$\ps$ for the brute-force check & \textbf{18} & qualitative only & 18\\
Liquid packing fraction & $\eta=0.526\pm0.005$ at $\ps=15$ (coexistence: $0.471\pm0.005$) & \yes & $\approx0.543$ at $\ps=18$ (validated parent)\\
$N$ (FFS) & 1000 & \yes & 1000\\
$N$ (umbrella sampling) & 8000 & \yes & not used\\
Dynamics & kinetic MC in NpT, small max.\ displacement & \yes & yes\\
dr for the initial flux $\phi_{l\to J}$ & $0.01\,\sigma$ & \yes & branch \texttt{dr\_0.01}\\
dr for FFS branching trajectories & $0.04\,\sigma$ & \yes & branch \texttt{dr\_0.04}\\
One MC cycle & one trial move per particle \textbf{plus one volume move} & \yes & 1000 + 1\\
Swap moves & none, or \textbf{20\%} & \yes & none implemented\\
Reaction coordinate & $n$ = particles in the largest solid cluster, from ref.\ [18] & citation only & see \texttt{04\_cluster\_size} (to be written)\\
\midrule
\multicolumn{4}{@{}l}{\textbf{\textcolor{red!70!black}{NOT stated in the paper}}}\\
\midrule
Volume step \texttt{vmax} & --- & \no & you tested it (\texttt{5\_Vmax\_Search})\\
dr of the brute-force $\ps=18$ run & --- & \no & you ran 0.01 and 0.04\\
System size at $\ps=18$ & --- & \no & 1000\\
FFS interface positions, trajectories per interface & --- & \no & n/a\\
Cluster thresholds (cutoff, dot-product, min.\ connections) & see ref.\ [18] & \no & q6 script: 1.45, $>0.70$, $\ge7$\\
MC cycles $\to$ physical time & --- & \no & n/a\\
Number of runs / seeds at $\ps=18$ & --- & \no & 1 per branch\\
\bottomrule
\end{longtable}
\normalsize

%--------------------------------------------------------------------
\section{The experiments, in order}

\begin{center}
\newcommand{\fb}[1]{\fcolorbox{accent}{soft}{\parbox{2.5cm}{\centering\small #1}}}
\fb{\textbf{Step 0}\\brute force\\$\ps=15$}\,$\rightarrow$\,\fb{\textbf{Step 1}\\FFS, no swaps\\$\ps=15$}\,$\rightarrow$\,\fb{\textbf{Step 2}\\FFS, 20\% swaps\\$\ps=15$}\\[6pt]
$\downarrow$\\[2pt]
\fb{\textbf{Step 3}\\brute force $\ps=18$\\(the one you ran)}\,$\rightarrow$\,\fb{\textbf{Step 4}\\umbrella sampling\\$N=8000$, $\ps=15$}
\end{center}

For each step: \textbf{why $\rightarrow$ what they did $\rightarrow$ how it was measured $\rightarrow$ what they saw $\rightarrow$ what it led to.}

\subsection{Step 0 \quad Plain brute force at $\ps=15$}
\begin{itemize}
\item \textbf{Why:} try the simplest thing first.
\item \textbf{Did:} ordinary NpT MC for $3\times10^{6}$ MC cycles at $\Ts=1$, $\ps=15$, starting from the metastable liquid.
\item \textbf{Saw:} \textbf{no spontaneous nucleation.}
\item \textbf{Led to:} nucleation at $\ps=15$ is too rare for brute force, so a rare-event method is needed.
\end{itemize}

\subsection{Step 1 \quad Nucleation rate by FFS, no swaps ($\ps=15$)}
\begin{itemize}
\item \textbf{Why:} get the real rate and see which crystal forms.
\item \textbf{Method} (paper Eqs.~3--4). The rate is a product of two factors:
\begin{equation*}
\phi_{l\to s}=\phi_{l\to J}\;P_{J\to s},\qquad
P_{J\to s}=P_{J\to K}\,P_{K\to L}\cdots P_{N\to s}.
\end{equation*}
\begin{enumerate}
\item $\phi_{l\to J}$: how often fluctuations in the liquid create a small crystalline cluster of $J=5$ particles. It comes from an ordinary long kinetic-MC run with \textbf{dr $=0.01\sigma$}.
\item $P_{J\to s}$: the probability that a 5-particle cluster grows to a bulk solid instead of dissolving. Each factor $P_{J\to K}$ is estimated by launching many short trajectories from stored configurations at cluster size $J$ and counting the fraction that reach size $K$ before the cluster dissolves. The successful ones become the starting points for the next interface. Trajectories from the same configuration differ only by the random seed, with \textbf{dr $=0.04\sigma$}. FFS only works when the dynamics are not deterministic; the seed makes them stochastic.
\end{enumerate}
\item \textbf{Reaction coordinate:} $n$ = number of particles in the largest solid cluster. It separates liquid from solid but is \emph{blind to the crystal structure}: every particle in disordered fcc \emph{or} CsCl counts as crystalline, while in the metastable liquid fewer than 5 of 1000 do.
\item \textbf{Saw} (numbers checked on the PDF, p.\,3):
\end{itemize}
\begin{center}
\begin{tabular}{@{}lll@{}}
\toprule
Quantity & Value & Note\\
\midrule
$P_{5\to s}$ & $10^{-28\pm2}$ & growth probability of a 5-particle cluster\\
$\phi_{l\to5}$ & $10^{-4\pm1}\,D_0\sigma^{-5}$ & $D_0$: diffusion coefficient at infinite dilution\\
Nucleation rate $\phi_{l\to s}$ & $10^{-32\pm3}\,D_0\sigma^{-5}$ & exponents add: $-4+(-28)=-32$\\
Critical nucleus $n^{*}$ & $120\pm15$ & size with 50\% chance to redissolve\\
\bottomrule
\end{tabular}
\end{center}
\begin{itemize}
\item \textbf{Structure of the growing cluster:} they compute $g(r)$ of \emph{only the particles in the cluster} and compare with bulk disordered-fcc and bulk CsCl curves (Fig.~2a, an 80-particle cluster). It matches \textbf{fcc-like}, even though the stable solid at $\Ts=1$ is CsCl (the paper writes ``$\ps=1$'', an apparent typo for 15).
\item \textbf{Led to:} the nucleus is the \emph{metastable} phase, not the stable one. This looks like Ostwald's step rule, which ST explains by ``the fcc barrier is lower''. That claim needs testing.
\end{itemize}

\subsection{Step 2 \quad Repeat FFS with 20\% swap moves ($\ps=15$)}
\begin{itemize}
\item \textbf{Why:} a direct test of the quasi-equilibrium assumption. If clusters are already in equilibrium, adding moves that only help them rearrange cannot change what they look like.
\item \textbf{Did:} same FFS, but the kinetic MC also has a \textbf{swap} trial move (exchange a $+$ and a $-$ particle). It is unphysical for colloids but removes the charge-ordering bottleneck.
\item \textbf{Saw:}
\end{itemize}
\begin{center}
\begin{tabular}{@{}lcc@{}}
\toprule
 & No swaps (Step 1) & 20\% swaps (Step 2)\\
\midrule
Cluster structure & disordered fcc-like & charge-ordered CsCl\\
Critical nucleus $n^{*}$ & $120\pm15$ & $65\pm15$\\
$P_{5\to s}$ & $10^{-28\pm2}$ & $10^{-15\pm1}$\\
\bottomrule
\end{tabular}
\end{center}
\begin{itemize}
\item \textbf{Led to:} the nucleation pathway \emph{changes} when sampling is improved, so \textbf{local equilibrium is not established in the natural (no-swap) dynamics.} That contradicts the key assumption of ST. Their explanation: the cluster grows too fast to explore its configurations and gets kinetically trapped in a metastable structure.
\end{itemize}

\subsection{Step 3 \quad Check at $\ps=18$ with no rare-event method (the step you reproduced)}
\begin{itemize}
\item \textbf{Why:} show that the effect is not an artifact of FFS.
\item \textbf{Did:} at $\ps=18$ the probability of forming a post-critical nucleus is higher, so no FFS is needed. They ran plain kinetic MC and a Brownian dynamics simulation.
\item \textbf{Saw:} kinetic MC \textbf{without swaps} forms a \textbf{substitutionally disordered fcc} lattice; Brownian dynamics gives the \textbf{same}; \textbf{with swaps} the liquid becomes a \textbf{substitutionally ordered} lattice.
\item \textbf{Led to:} the same phenomenon, now seen directly. It also supports MC as a stand-in for real colloid dynamics.
\end{itemize}
\begin{warnbox}
\textbf{No numbers are given for this step:} no cluster sizes, times, $N$, dr or vmax. Your target here is therefore \emph{qualitative}.
\end{warnbox}

\subsection{Step 4 \quad Free-energy barriers by umbrella sampling ($\ps=15$, $N=8000$)}
\begin{itemize}
\item \textbf{Why:} Steps 1--3 show which crystal forms. To contradict ST you must also show that the \emph{other} crystal has the lower barrier.
\item \textbf{Method:} umbrella sampling holds the system near each value of $n$ in turn and assembles the free energy $\Delta G(n)$. (The paper only says it ``determines the variation of the free energy with a reaction coordinate''; the bias and the reweighting are not described.)
\item \textbf{Done twice:}
\begin{enumerate}
\item \emph{No swaps:} CsCl clusters grow very slowly; fcc clusters form and persist. If the run is long enough, fcc clusters eventually turn into CsCl.
\item \emph{With swaps:} clusters grow directly in their lowest-free-energy state (CsCl).
\end{enumerate}
\item \textbf{Saw} (Fig.~3a, \emph{approx.}\ read from the plot):
\end{itemize}
\begin{center}
\begin{tabular}{@{}lcc@{}}
\toprule
 & CsCl & Disordered fcc\\
\midrule
Barrier $\Delta G^{*}/\kT$ & $\approx20$--$22$ & $\approx31$--$32$\\
$n$ at the maximum & $\approx60$--$100$ & $\approx120$--$140$\\
\bottomrule
\end{tabular}
\end{center}
The disordered critical cluster is \emph{bigger} than the ordered one (snapshots, Fig.~3b,c).
The authors warn that a barrier height can depend on the choice of reaction coordinate, whereas the rate is a true observable.

%--------------------------------------------------------------------
\section{How the argument is stitched together}

\begin{enumerate}
\item \textbf{Step 1:} without swaps, \textbf{fcc-like} nuclei form (rate $\approx10^{-32}D_0\sigma^{-5}$).
\item \textbf{Step 4:} the \textbf{CsCl} nucleus has the \textbf{lower} barrier. So the system nucleates the phase with the \emph{higher} barrier. That alone already contradicts ST.
\item \textbf{Step 2:} with swaps the nucleus is CsCl and nucleation is far easier. Therefore the no-swap system was \emph{not} exploring all structures during nucleation: it is \textbf{out of equilibrium}.
\item \textbf{Step 3:} plain MC and Brownian dynamics at $\ps=18$ show the same, so it is not an FFS artifact.
\item \textbf{Mechanism (their interpretation), ``self-poisoning'':} small fcc clusters grow so fast that they never become CsCl. An fcc cluster would have to nearly redissolve before a CsCl one could form, so the early success of fcc blocks the more stable CsCl.
\item \textbf{Why CNT is not rescued by $\Gamma$:} attachment and detachment rates are about the same for fcc and CsCl clusters, so a large prefactor cannot compensate for the higher barrier.
\end{enumerate}

\begin{keybox}
\textbf{Headline:} the phase that crystallizes first is not the one with the lowest barrier, because sub-critical nuclei are not in quasi-equilibrium.
\end{keybox}

%--------------------------------------------------------------------
\section{Measurement techniques in detail}

\subsection{Cluster size $n$}
``The number of particles in the largest solid cluster'', citing ref.\ [18]. It is structure-blind (fcc and CsCl particles both count).
The thresholds are not in this paper. Your q6 script implements a bond-orientational version of this idea with its own parameter values. The comparison with the reference is the subject of the next note (\texttt{04\_cluster\_size\_comparison}, not yet written; it needs the ref.\ [18] PDF).

\subsection{Critical nucleus}
The size at which a cluster has a 50\% chance to shrink and 50\% to grow. Measured by launching trajectories from clusters of a given size and counting how many redissolve: $120\pm15$ without swaps, $65\pm15$ with 20\% swaps.

\subsection{FFS rate}
Rate $=$ (flux to the first interface) $\times$ (product of interface-to-interface probabilities). The errors are quoted on the exponent ($\pm2$ on the probability, $\pm3$ on the rate), so the numbers are good to orders of magnitude only.

\subsection{Crystal structure from $g(r)$}
$g(r)$ of the cluster particles only, compared with the two bulk curves (the cluster curve is multiplied by 20 to be visible).
\begin{itemize}
\item A small cluster is mostly surface, so its $g(r)$ is noisy.
\item The total $g(r)$ cannot separate \emph{disordered} from \emph{ordered} fcc, nor fcc from hcp. The word ``disordered'' rests on the dynamics (without swaps charges cannot rearrange), not on a measurement shown in the paper.
\end{itemize}

\subsection{Free-energy barrier}
The $\Delta G(n)$ curve from umbrella sampling; the barrier is its maximum (Fig.~3a).

\subsection{Time and rate units}
Rates are in $D_0\sigma^{-5}$ (per volume per time, hydrodynamics ignored). There is no calibration of MC cycles to physical time, so cycles at different dr are not comparable.

%--------------------------------------------------------------------
\section{Map to your work}

\begin{center}
\begin{tabularx}{\textwidth}{@{}>{\raggedright\arraybackslash}p{7.6cm} X@{}}
\toprule
\textbf{Paper step} & \textbf{In your project}\\
\midrule
Model, NpT MC, $N=1000$, $\Ts=1$, $\kappa\sigma=6$, $r_c=3.5\sigma$, 1 volume move per cycle & \done\\
Step 3, no-swap half (brute force $\ps=18$) & \done: two branches, dr $=0.01$ and $0.04$, $3\times10^{6}$ cycles each; only dr $=0.04$ nucleated\\
Step 3, swap half; Brownian dynamics & \notdone{} (no swap move in the code)\\
Step 0 ($\ps=15$ brute force) & \notdone\\
Steps 1, 2 (FFS) & \notdone{} (deliberate choice)\\
Step 4 (umbrella sampling) & \notdone\\
\bottomrule
\end{tabularx}
\end{center}

%--------------------------------------------------------------------
\section{If you redo it: checklist}

\begin{enumerate}
\item Prepare a liquid at $\Ts=1$, $\ps=18$. Your validated parent is \texttt{4\_Clean\_NoSwap\_Replication/\allowbreak common/\allowbreak validated\_parent\_\allowbreak N1000\_T1\_P18.res} ($\eta\approx0.543$).
\item Run kinetic NpT MC with 1000 displacement moves $+$ 1 volume move per cycle, no swaps, at the dr you choose. The paper gives 0.01 and 0.04 for specific roles but not for this run.
\item Choose vmax. The paper does not state it; it changes the dynamics, not the NpT equilibrium.
\item Compute the cluster size with a \emph{structure-blind} criterion (see \texttt{04\_...} once written). A PTM ``FCC/HCP only'' count excludes bcc and would miss a CsCl nucleus.
\item A nucleation event looks like $n$ growing steadily far beyond about 100 particles and never returning to near zero.
\item Check the nucleus structure: total $g(r)$ as in Fig.~2, and ideally $g(r)$ split by charge to see whether charges are random (disordered fcc) or ordered (CsCl).
\item Repeat with more seeds if you want any statistics. The paper gives none for $\ps=18$.
\end{enumerate}

%--------------------------------------------------------------------
\section{Glossary}

\small
\begin{center}
\begin{tabularx}{\textwidth}{@{}>{\bfseries}l X@{}}
\toprule
Term & Meaning\\
\midrule
NpT MC & Monte Carlo at constant particle number, pressure, temperature (moves: particle displacement and volume change)\\
MC cycle & here: 1000 displacement attempts $+$ 1 volume attempt\\
dr & maximum displacement per particle move (in $\sigma$); sets the speed of the dynamics\\
vmax & maximum change of $\ln V$ in a volume move; an algorithmic parameter\\
ST & Stranski--Totomanow conjecture (the lowest-barrier phase nucleates)\\
CNT & classical nucleation theory, $R=\Gamma e^{-\Delta G_{\rm crit}/\kT}$\\
FFS & forward-flux sampling, a rare-event rate method\\
Umbrella sampling & biased sampling used to compute free-energy profiles\\
Disordered fcc & positions on an fcc lattice, charges assigned randomly\\
CsCl & positions on a bcc lattice with alternating charges (ordered)\\
Swap move & exchange of a $+$ and a $-$ particle; unphysical for colloids\\
Ostwald step rule & a metastable phase forms first and later transforms to the stable one\\
q6, PTM, CNA & structure-detection tools used in the simulations; see \texttt{04\_...}\\
\bottomrule
\end{tabularx}
\end{center}

\end{document}
